\documentclass[preview]{standalone}

\usepackage{amsmath}
\usepackage{amssymb}
\usepackage{stellar}
\usepackage{definitions}

\begin{document}

\id{group-isomorphism-theorems}
\genpage

\section{Isomorphism theorems}

\begin{plainhtml}
    The same theorems apply to rings, vector spaces and such.
\end{plainhtml}

\begin{snippettheorem}{group-first-isomorphism-theorem}{First isomorphism theorem}
    Let \(\varphi\colon G \fromto H\) be a \grouphomomorphism.
    Then, \[G / \grpker \varphi \groupisomorphic \image \varphi\]
\end{snippettheorem}

\begin{snippettheorem}{group-second-isomorphism-theorem}{Second isomorphism theorem}
    Let \(G\) be a \group and let
    \(H \subgroupleq G, K \normalsubgrp G\).
    Then,
    \[
        H \intersection K \normalsubgrp H
        \land \frac{H}{H \intersection K} \groupisomorphic \frac{HK}{K}
    \]
\end{snippettheorem}

\begin{snippettheorem}{group-third-isomorphism-theorem}{Third isomorphism theorem}
    Let \(G\) be a \group and let \(N\subgroupleq H\normalsubgrp G\), with
    \(N\normalsubgrp G\). Then \(H/N\normalsubgrp G/N\) and
    \[
        \frac{G/N}{H/N} \groupisomorphic \frac{G}{H}
    \]
\end{snippettheorem}

\begin{snippet}{isomorphism-theorems-exp11}
    Let \(\varphi\colon G\fromto H\) be a \grouphomomorphism. For
    \(K\subgroupleq G\) and \(L\subgroupleq H\),
    \begin{enumerate}
        \item \(K\subgroupleq\varphi^{-1}(\varphi(K))\), with equality \ifandonlyif
            \(\grpker\varphi\subgroupleq K\);
        \item \(\varphi(\varphi^{-1}(L))=L\intersection\image\varphi\), so equality
            with \(L\) holds \ifandonlyif \(L\subgroupleq\image\varphi\).
    \end{enumerate}
\end{snippet}

\begin{snippettheorem}{group-subgroup-correspondence-theorem}{Subgroup correspondence theorem}
    Let \(G\) be a \group and \(N\normalsubgrp G\). The assignments
    \[
        H\longmapsto H/N,
        \qquad
        K\longmapsto\pi_N^{-1}(K)
    \]
    give mutually inverse, inclusion-preserving bijections between the
    \subgroup[subgroups] \(H\subgroupleq G\) containing \(N\) and the
    \subgroup[subgroups] \(K\subgroupleq G/N\). Under this bijection,
    \[
        H\normalsubgrp G \iff H/N\normalsubgrp G/N,
        \qquad |G:H|=|G/N:H/N|.
    \]
\end{snippettheorem}

\begin{snippetproof}{group-subgroup-correspondence-theorem-proof}{group-subgroup-correspondence-theorem}{Subgroup correspondence theorem}
    Apply the image--preimage identities to the surjective canonical
    projection \(\pi_N\colon G\fromto G/N\), whose kernel is \(N\).
    They show that \(H\mapsto\pi_N(H)=H/N\) and inverse image are mutually
    inverse on the stated families. Images and preimages under a
    homomorphism preserve normality in the relevant ambient image, giving
    the normality equivalence. Finally, the cosets of \(H\) in \(G\)
    correspond bijectively to the cosets of \(H/N\) in \(G/N\).
\end{snippetproof}

\subsection{Proofs}

\begin{snippetproof}{group-first-isomorphism-theorem-proof}{group-first-isomorphism-theorem}{First isomorphism theorem}
    Let \(N=\grpker\varphi\). By the universal property of \(G/N\),
    \(\varphi\) induces a homomorphism
    \[
        \overline{\varphi}\colon G/N\fromto\image\varphi,
        \qquad \overline{\varphi}(gN)=\varphi(g).
    \]
    It is surjective by the definition of \(\image\varphi\). If
    \(\overline{\varphi}(gN)=1_H\), then \(g\in\grpker\varphi=N\), so
    \(gN=N\). Thus its kernel is trivial, and it is injective. Therefore it
    is an isomorphism.
\end{snippetproof}

\begin{snippetproof}{group-second-isomorphism-theorem-proof}{group-second-isomorphism-theorem}{Second isomorphism theorem}
    Since \(K\normalsubgrp G\), the product \(HK\) is a subgroup of \(G\)
    and \(K\normalsubgrp HK\). Consider
    \[
        \varphi\colon H\fromto HK/K,
        \qquad \varphi(h)=hK.
    \]
    It is surjective because every coset in \(HK/K\) has the form
    \(hkK=hK\). Moreover,
    \[
        \grpker\varphi=\{h\in H\suchthat hK=K\}=H\intersection K.
    \]
    Hence \(H\intersection K\normalsubgrp H\), and the first isomorphism
    theorem gives \(H/(H\intersection K)\groupisomorphic HK/K\).
\end{snippetproof}

\begin{snippetproof}{group-third-isomorphism-theorem-proof}{group-third-isomorphism-theorem}{Third isomorphism theorem}
    The subgroup correspondence theorem gives
    \(H/N\normalsubgrp G/N\). Define
    \[
        \theta\colon G/N\fromto G/H,
        \qquad \theta(gN)=gH.
    \]
    This is well-defined because \(N\subgroupleq H\), and it is a surjective
    homomorphism. Its kernel is \(H/N\). The first isomorphism theorem yields
    \[
        \frac{G/N}{H/N}\groupisomorphic G/H.
    \]
\end{snippetproof}

\end{document}
